\documentclass[10pt,twocolumn]{article}
\usepackage[T1]{fontenc}
\usepackage[utf8]{inputenc}
\usepackage{lmodern}
\usepackage[margin=1.8cm,top=2.4cm,bottom=2cm,columnsep=0.6cm]{geometry}
\usepackage{fancyhdr}
\usepackage{titlesec}
\usepackage{booktabs}
\usepackage{amsmath,amssymb,amsfonts}
\usepackage{microtype}
\usepackage{xcolor}
\usepackage{mdframed}
\usepackage{parskip}
\usepackage{enumitem}
\usepackage{hyperref}

\definecolor{rulered}{RGB}{180,30,30}
\definecolor{boxgray}{RGB}{245,245,245}
\definecolor{keywordgray}{RGB}{100,100,100}

\pagestyle{fancy}
\fancyhf{}
\fancyhead[L]{\textsc{\small Claude Mag}}
\fancyhead[C]{\small\textit{Machine Learning Research Digest}}
\fancyhead[R]{\small 2026-09-30}
\fancyfoot[C]{\small\thepage}
\renewcommand{\headrulewidth}{0.8pt}

\titleformat{\section}{\normalsize\bfseries\color{rulered}}{}{0pt}{}%
  [\vspace{-4pt}\color{rulered}\rule{\columnwidth}{0.4pt}]
\titlespacing{\section}{0pt}{8pt}{4pt}

\hypersetup{colorlinks=true,urlcolor=rulered,linkcolor=black}

\begin{document}
\setlength{\parindent}{0pt}
\setlength{\parskip}{4pt}

{\small\color{keywordgray}\textbf{Keywords:}
  visual token pruning \textbullet{}
  multimodal LLMs \textbullet{}
  conditional removability \textbullet{}
  inference efficiency}\par\vspace{4pt}

{\LARGE\bfseries\raggedright
  CoRePrune}\par\vspace{4pt}

{\large\itshape\raggedright
  From Token Importance to Conditional Removability:
  Rethinking Visual Token Pruning in Multimodal Large Language Models}\par\vspace{6pt}

{\small\textbf{By} Shengli He et al.}\par\vspace{2pt}

{\small\color{keywordgray}arXiv:2609.26484 \textbullet{} September 2026}
\par\vspace{2pt}\noindent{\color{rulered}\rule{\columnwidth}{0.6pt}}\vspace{6pt}

Token importance and redundancy are poor proxies for removability because
they ignore representation depth and deletion-set context; CoRePrune's
progressive depth-aware perturbation tracking and set-conditioned refinement
retain 90.3\% of dense-model performance with 128 visual tokens while cutting
prefill time by 51\% across five multimodal LLM backbones.

\begin{mdframed}[backgroundcolor=boxgray,linecolor=rulered,linewidth=1pt,
    innerleftmargin=6pt,innerrightmargin=6pt,
    innertopmargin=6pt,innerbottommargin=6pt]
\textbf{\small AT A GLANCE}\par\vspace{4pt}
\begin{description}[leftmargin=0pt,labelsep=4pt,itemsep=2pt,parsep=0pt,
    font=\small\bfseries]
  \item[Problem:] Visual token pruning methods score tokens by importance or
    redundancy but miss that removability is conditioned on transformer depth
    and the surrounding deletion set.
  \item[Why it matters:] Incorrect removability estimates cause unexpected
    quality drops at aggressive budgets; fixing the estimate recovers performance.
  \item[Method:] Two-stage CoRePrune: progressive perturbation-aware visual
    pruning across depth, followed by set-conditioned refinement after
    visual-text interaction.
  \item[Result:] 90.3\% dense-model performance with 128 tokens; 51.0\%
    prefill speedup on five MLLM backbones.
  \item[Evidence:] Standard image, high-resolution, and video benchmarks
    across LLaVA-1.5, InternVL2, LLaVA-HR, and VideoLLM.
\end{description}
\end{mdframed}

\section{The Big Picture}

Multimodal large language models (MLLMs) convert images into sequences of
visual tokens---typically hundreds to thousands---that dominate the cost of
the prefill phase. Training-free token pruning addresses this by scoring each
visual token and discarding the lowest-scoring ones before or during the
forward pass. The appeal is obvious: no fine-tuning, no architecture change,
just a cheaper forward pass.

The trouble is that existing scoring functions---attention weight, activation
norm, pairwise similarity---treat removability as a static per-token property.
CoRePrune shows that it is not. Removing the same token at layer 5 versus
layer 20 causes substantially different output perturbations because the
representations have been increasingly fused with context. Removing a token
when 100 others are already gone has a different effect than removing it from
the full set because the remaining tokens have compensated for prior deletions.
These two dependencies---\emph{depth} and \emph{deletion set}---are precisely
what existing methods ignore.

\section{Why Existing Approaches Fall Short}

\textbf{Importance-based methods} (attention score, gradient norm) score tokens
in isolation and at a single layer; they cannot predict how importance changes
as representations propagate through depth.

\textbf{Redundancy-based methods} (cosine similarity) identify tokens that are
similar to their neighbours, but spatial similarity at one layer does not imply
similarity after additional self-attention mixing.

\textbf{Static depth assignment:} Most methods apply pruning at a fixed layer
(early or late), ignoring that the same token may be safe to remove at depth 20
but harmful to remove at depth 5.

\textbf{Set-agnostic scoring:} All existing methods score each candidate
independently, ignoring that removing token $v_i$ with $v_j$ already gone
may have a very different effect than removing $v_i$ from the full set.

\section{Core Mechanism}

\textbf{Conditional Removability.} CoRePrune operationalises removability as
the downstream perturbation caused by removing a token or set of tokens,
measured at the current depth under the current deletion set. This quantity
is tracked dynamically rather than estimated once at a fixed layer.

\textbf{Stage 1: Progressive Perturbation-Aware Visual Pruning (PPVP).}
As the forward pass proceeds through successive transformer layer blocks,
PPVP re-estimates the perturbation each candidate removal would cause. A token
whose removal perturbation \emph{shrinks} as depth increases has been
compensated for by the model's own representation mixing---it is safe to remove.
Tokens whose perturbation grows are promoted in priority.

\textbf{Stage 2: Set-Conditioned Refinement (SCR).}
After the first visual-text interaction layer (where visual and text tokens
first attend to each other), SCR re-evaluates each remaining candidate under
the \emph{current} deletion set. Tokens that now provide net benefit given
what has already been removed are recovered; tokens whose marginal contribution
has collapsed are added to the deletion set.

\section{Technical Formulation}

Let $V = \{v_1, \ldots, v_N\}$ be the visual token sequence and
$f_l(S)$ the representations at layer $l$ for token set $S$.

\textbf{Perturbation at depth $l$ for deletion set $D$:}
\[
P_l(D) = \bigl\|f_l(V \setminus D) - f_l(V)\bigr\|_F
\]

\textbf{PPVP priority update for candidate $v_i$:}
\[
\text{safe at } l_2 \iff P_{l_2}(\{v_i\}) < P_{l_1}(\{v_i\}),
\quad l_2 > l_1
\]

\textbf{SCR marginal contribution} under current deletion set $D$ at layer $l^*$:
\[
r(v_j; D) = P_{l^*}(D) - P_{l^*}(D \cup \{v_j\})
\]
Recover $v_j$ if $r(v_j; D) > 0$;
add $v_j$ to $D$ if $r(v_j; D) \ll 0$.

\section{Learning or Inference Procedure}

CoRePrune is \textbf{training-free} and runs during the prefill phase:
\begin{enumerate}[itemsep=2pt,parsep=0pt]
  \item Forward pass to layer $l_1$: compute initial perturbation estimates.
  \item Continue to layer $l_2 > l_1$: update estimates via PPVP; rank candidates.
  \item Tentatively remove lowest-priority tokens to meet the token budget.
  \item Continue to visual-text interaction layer $l^*$.
  \item Apply SCR: recover tokens with $r(v_j; D) > 0$; finalise deletion set.
  \item Continue forward pass with the pruned set; generate output.
\end{enumerate}

Overhead: two additional perturbation probes during prefill, negligible relative
to the attention savings from the reduced token count.

\section{What the Guarantee Says}

PPVP provides a constructive argument: if $P_{l_2}(\{v_i\}) < P_{l_1}(\{v_i\})$,
the model has already compensated for $v_i$'s absence by depth $l_2$, making
removal safe. SCR provides a greedy local-optimum guarantee: after SCR, no
single token swap (recovery or addition) would reduce the deletion-set perturbation.

\section{Experimental Findings}

\begin{itemize}[itemsep=2pt,parsep=0pt]
  \item \textbf{5 MLLM backbones:} LLaVA-1.5, InternVL2, LLaVA-HR,
    VideoLLM, and one additional model
  \item At \textbf{128 visual tokens}: CoRePrune retains \textbf{90.3\%}
    of dense-model performance (vs.\ 83--86\% for prior methods) and reduces
    aggregate prefill time by \textbf{51.0\%} (vs.\ 45--48\% prior)
  \item At 64 tokens: advantage over baselines grows due to stronger
    deletion-set interaction effects
  \item Largest gains on global reasoning tasks (VQA, video QA); smaller
    gains on local patch tasks where interaction effects are weaker
\end{itemize}

\section{Ablations and Interpretation}

\textbf{PPVP alone} removes the depth-static assumption and recovers half
the total gain over attention-based baselines.

\textbf{SCR alone} removes the set-agnostic assumption and recovers the
other half. The joint gain exceeds the sum, revealing non-linear interaction.

\textbf{PPVP probe depth:} Tracking across $\geq$3 layer blocks is necessary;
fewer blocks converges to depth-static behaviour.

\textbf{SCR timing:} Applying SCR before visual-text interaction or at the
final layer both degrade performance; the visual-text fusion layer is the
optimal intervention point.

\section*{Reference}

Shengli He et al.
\textbf{From Token Importance to Conditional Removability: Rethinking
Visual Token Pruning in Multimodal Large Language Models}.
arXiv:2609.26484, September 2026.\\
\url{https://arxiv.org/abs/2609.26484}

\end{document}
